\documentclass[12pt]{article}

\usepackage[spanish]{babel}
\usepackage[utf8]{inputenc}
\usepackage[T1]{fontenc}
\usepackage{amsmath,amssymb}
\usepackage{geometry}
\usepackage{array}
\usepackage{enumitem}
\usepackage{lmodern}

\geometry{letterpaper, margin=2cm}

\begin{document}
	
	\begin{center}
		{\LARGE \textbf{Autoevaluación Evaluación II}}\\[0.2cm]
		{\large Contextualizada en Agronomía}
	\end{center}
	
	\vspace{0.4cm}
	
	\textbf{Instrucciones:} En cada caso realice todo el desarrollo en forma ordenada y no olvide responder la pregunta planteada e interpretar el resultado en contexto agrícola.
	
	\vspace{0.5cm}
	
	\begin{enumerate}
		
		%-------------------------------------------------
		\item \textbf{Registro de problemas en predios agrícolas}
		
		Un agrónomo analiza los problemas más comunes que afectan distintos predios agrícolas.  
		Sea \(x\) el número base de predios monitoreados en una zona agrícola.
		
		Se obtiene la siguiente información:
		
		\begin{center}
			\begin{tabular}{|m{8cm}|>{\centering\arraybackslash}m{4cm}|}
				\hline
				\textbf{Problemas en el cultivo} & \textbf{Número de predios} \\
				\hline
				Daño por plagas & \(x + 5\) \\ \hline
				Enfermedades fúngicas & \(2x - 3\) \\ \hline
				Déficit hídrico & \(x - 2\) \\ \hline
				Daño por heladas & \(x - 7\) \\ \hline
			\end{tabular}
		\end{center}
		
		\begin{enumerate}[label=\alph*)]
			\item Determine el número total de predios afectados por estos problemas.
			\item Determine la diferencia entre los predios afectados por plagas y por heladas.
		\end{enumerate}
		
		%-------------------------------------------------
		\item \textbf{Aplicación de fertilizante}
		
		Para mejorar el rendimiento de un cultivo se deben fertilizar:
		\[
		(3n^2+5n)
		\]
		plantas y la dosis a aplicar a cada una es de
		\[
		(2n+7)\text{ mL}.
		\]
		
		Determine la cantidad total de fertilizante (en mL) que se debe aplicar.
		
		%-------------------------------------------------
		\item \textbf{Tratamiento fitosanitario}
		
		En un cultivo afectado por una plaga se prescribe un tratamiento. Para el tratamiento completo se requieren \(x\) envases del producto, cada uno con:
		\[
		(x^2-x+72)\text{ mL}
		\]
		
		Si el cultivo requiere una aplicación diaria de:
		\[
		(x^2-9x)\text{ mL}
		\]
		
		¿Cuántos días dura el tratamiento?
		
		%-------------------------------------------------
		\item \textbf{Costo de insumos agrícolas}
		
		Un agricultor adquiere insumos en una cantidad dada por:
		\[
		Q=\frac{x^2-5x+6}{x-3}
		\]
		
		a un precio por unidad de:
		\[
		P=\frac{x^3+x^2+4x+4}{x^2-x-2}
		\]
		
		El costo total es:
		\[
		C=P\cdot Q
		\]
		
		\begin{enumerate}[label=\alph*)]
			\item Simplifique el costo total.
			\item Evalúe el costo para \(x=1500\).
		\end{enumerate}
		
		%-------------------------------------------------
		\item \textbf{Modelo de concentración de agroquímicos}
		
		La concentración de un agroquímico en el suelo está dada por:
		
		\[
		C=\left[
		\frac{x^2-x-12}{x^2-9}
		\cdot
		\frac{x^3+x^2-2x}{x^2-2x-8}
		\div
		\frac{x-1}{2x^2-6x}
		\right]
		\]
		
		Simplifique la expresión aplicando factorización.
		
		\vspace{0.3cm}
		
		Además, simplifique:
		
		\[
		\frac{\bigl[(8x^2-5)+4\bigr]^2(2x+7)^2-(4x-7)(3x-2)(3x+4)}
		{(5x+3)(3+5x)-9}
		\]
		
	\end{enumerate}
	
\end{document}